\chapter{Session keys}
\label{ch:sesskeys}

\chapterbrief
  {Derive the two per-transaction AES keys from the nonces exchanged during authentication.}
  {Private static helpers, and enums instead of boolean flags.}
  {EV2 key derivation verified against independent vectors.}

\section{The problem: never use the long-term key directly}

Your card holds a permanent AES-128 key. It would be technically possible to
encrypt every message with it. It would also be a bad idea, for three reasons
worth understanding:

\begin{description}[leftmargin=0pt, style=nextline]
  \item[Replay.] If one key encrypts everything forever, a recorded ``debit
        1.50'' frame stays valid. An attacker records one purchase and replays it
        --- or replays a \emph{credit}.
  \item[Traffic analysis.] The same plaintext under the same key gives the same
        ciphertext every time. An eavesdropper learns which commands are which
        without breaking anything.
  \item[Exposure.] Every message under the long-term key is another sample for
        cryptanalysis. Keys wear out with use.
\end{description}

So DESFire does what TLS does: use the long-term key \emph{once}, to bootstrap a
pair of throwaway keys valid only for this tap.

\begin{analogy}{The hotel key card}
A hotel does not hand guests the master key that opens every door forever. You
present identification at the desk once --- the long-term secret --- and they
programme a card that opens \emph{your} room, \emph{this} week. Lose it and the
damage is bounded; the master key never left the safe.

The nonces are what make each session unique. Both sides contribute randomness,
so neither can force a repeat --- which matters, because if the card alone chose,
a compromised card could replay the same session forever.
\end{analogy}

\textbf{Two keys, not one}, because encryption and authentication are different
jobs and reusing one key for both is a classic footgun:

\begin{itemize}
  \item $K_{enc}$ --- encrypts data. \emph{Confidentiality}: nobody can read the balance.
  \item $K_{mac}$ --- authenticates commands. \emph{Integrity}: nobody can alter the debit.
\end{itemize}

\section{The derivation}

Both keys come from the same construction with a different label byte. The input
is a 32-byte ``session vector'' (SV):

\begin{center}
\begin{tikzpicture}[font=\scriptsize\sffamily, x=4.6mm, y=8mm]
  \definecolor{c1}{HTML}{DCE6F2}
  \definecolor{c2}{HTML}{F2E4CC}
  \definecolor{c3}{HTML}{D8ECE0}
  \definecolor{c4}{HTML}{EFE0EC}
  % row of 32 cells
  \foreach \i in {0,...,31}{ \draw[brewergrey!60] (\i,0) rectangle ++(1,1); }
  \fill[c1] (0,0) rectangle (2,1);   \node at (1,0.5) {\tiny lbl};
  \fill[black!5] (2,0) rectangle (6,1); \node at (4,0.5) {\tiny 00 01 00 80};
  \fill[c2] (6,0) rectangle (8,1);   \node at (7,0.5) {\tiny A$_{0,1}$};
  \fill[c3] (8,0) rectangle (14,1);  \node at (11,0.5) {\tiny A$_{2..7}\oplus$B$_{0..5}$};
  \fill[c2] (14,0) rectangle (24,1); \node at (19,0.5) {\tiny rndB$_{6..15}$};
  \fill[c4] (24,0) rectangle (32,1); \node at (28,0.5) {\tiny rndA$_{8..15}$};
  \foreach \i in {0,...,31}{ \draw[brewergrey!60] (\i,0) rectangle ++(1,1); }
  % byte index ruler
  \foreach \i/\lbl in {0/0, 6/6, 8/8, 14/14, 24/24, 32/32}{
     \node[color=brewergrey, font=\tiny] at (\i,-0.35) {\lbl};
     \draw[brewergrey!50] (\i,0) -- (\i,-0.18);
  }
  \node[color=brewergrey, font=\tiny, anchor=west] at (0,1.45)
       {32-byte session vector \quad$\rightarrow$\quad session key $=$ CMAC(authKey, SV)};
\end{tikzpicture}
\end{center}

\begin{trace}
SV[0]      label            0xA5 for Kenc, 0x5A for Kmac
SV[1]      counter-label    the other one (0x5A when label is 0xA5)
SV[2..5]   00 01 00 80      fixed
SV[6..7]   rndA[0..1]
SV[8..13]  rndA[2..7] XOR rndB[0..5]      <- the interleave
SV[14..23] rndB[6..15]
SV[24..31] rndA[8..15]

session key = CMAC(authKey, SV)           32 bytes in, 16 out
\end{trace}

The CMAC you built in Chapter~\ref{ch:cmac} is doing duty here as a \textbf{key
derivation function}: it mixes a lot of input down into sixteen bytes in a way
that cannot be inverted. That is the same property that makes it a good MAC.

\begin{deeper}{Why the strange interleave?}
Both sides must derive identical keys, so the layout is ultimately just a
specification both implement. But note \emph{why it is mixed rather than
concatenated}: bytes from both nonces are XOR'd together, so neither party can
steer the result by choosing their nonce cleverly.

If the layout were plain \code{rndA || rndB}, a malicious reader could pick
\code{rndA} to force structure into the derived key. The interleave denies both
sides that control.

And \hex{0xA5}/\hex{0x5A} are bit-complements of each other, chosen to be
maximally distinct so the two derivations cannot accidentally collide. The label
sits at byte 0, so the very first AES block of the CMAC already differs and the
two computations diverge immediately.
\end{deeper}

\section{Build it}

\begin{jcode}
package com.midnightbrewer.hardware.crypto;

import java.util.Arrays;

/**
 * One authenticated session with a card: the two derived keys, and (from
 * Chapter 10) the secure-messaging state that goes with them.
 *
 * The keys are derived ONCE in the constructor and never change -- the same
 * invariant discipline as AesCmac. If a DesfireSession exists, its keys are
 * valid and belong together.
 */
public class DesfireSession {

    private final byte[] keyEnc;
    private final byte[] keyMac;

    /**
     * Which key we are deriving.
     *
     * An ENUM, not a boolean. deriveSessionKey(key, a, b, true) at a call site
     * tells you nothing -- true what? SessionKeyType.ENC is self-documenting,
     * the compiler rejects anything else, and adding a third key type later
     * does not change any existing signature.
     */
    private enum SessionKeyType { ENC, MAC }

    public DesfireSession(byte[] authKey, byte[] rndA, byte[] rndB) {
        if (authKey.length != 16) {
            throw new IllegalArgumentException("auth key must be 16 bytes");
        }
        this.keyEnc = deriveSessionKey(authKey, rndA, rndB, SessionKeyType.ENC);
        this.keyMac = deriveSessionKey(authKey, rndA, rndB, SessionKeyType.MAC);
    }

    // =================================================================
    // private static: it is a pure function (same inputs -> same output,
    // touches no instance state), so it does not need 'this'. Marking it
    // static says exactly that, and lets the constructor call it before the
    // object is finished being built.
    // =================================================================
    private static byte[] deriveSessionKey(byte[] authKey, byte[] rndA,
                                           byte[] rndB, SessionKeyType type) {
        if (rndA.length != 16 || rndB.length != 16) {
            throw new IllegalArgumentException("nonces must be 16 bytes");
        }
        byte[] sv = new byte[32];

        // TODO 1: sv[0] = 0xA5 for ENC, 0x5A for MAC.
        //         sv[1] = the opposite one.
        // TODO 2: sv[2..5] = 0x00, 0x01, 0x00, 0x80
        // TODO 3: sv[6] = rndA[0];  sv[7] = rndA[1];
        // TODO 4: sv[8..13] = rndA[2..7] XOR rndB[0..5]
        //         (Arrays.copyOfRange + Aes.xor is the tidy way)
        // TODO 5: sv[14..23] = rndB[6..15]     (10 bytes)
        // TODO 6: sv[24..31] = rndA[8..15]     (8 bytes)
        // TODO 7: return new AesCmac(authKey).compute(sv);
        return new byte[16];
    }

    /** Defensive copies -- callers must not be able to scribble on our keys. */
    public byte[] encKey() { return keyEnc.clone(); }

    public byte[] macKey() { return keyMac.clone(); }
}
\end{jcode}

\section{Two design notes worth absorbing}

\subsection{Enum over boolean}
\index{enum}

Compare these call sites:

\begin{jcode}
deriveSessionKey(key, rndA, rndB, true);                 // true... what?
deriveSessionKey(key, rndA, rndB, SessionKeyType.ENC);   // obvious
\end{jcode}

A boolean parameter is a note-to-self that the reader does not get to see. This
is sometimes called the \textbf{boolean trap}. Any time you are about to add a
\code{boolean} parameter that selects \emph{behaviour} rather than expressing
\emph{a yes/no fact}, an enum is better: self-documenting at the call site,
extensible without changing signatures, and impossible to pass the wrong way
round.

\subsection{Do not keep scratch work as state}

A tempting first draft stores the SV and the derived key as fields:

\begin{jcode}
private byte[] SV;      // don't
private byte[] SK;      // don't
\end{jcode}

Now the object has fields that are only meaningful \emph{during} construction.
Anyone reading the class must work out whether they are still valid later, and a
second call to a derive method would silently clobber them.

\begin{concept}{Fields versus locals}
\textbf{Scratch work belongs in local variables. Fields are for state the object
genuinely \emph{has}.}

\code{keyEnc} and \code{keyMac} are state. \code{sv} is a temporary. The test:
\emph{``if someone asked this object about this field a minute from now, would
the answer be meaningful?''} If not, it is a local.
\end{concept}

\section{Verify}

These vectors were generated independently with Python's \code{cryptography}
library, so agreement means your Java is \emph{right} --- not merely
self-consistent.

\begin{trace}
authKey = 00112233445566778899aabbccddeeff
rndA    = a0a1a2a3a4a5a6a7a8a9aaabacadaeaf
rndB    = b0b1b2b3b4b5b6b7b8b9babbbcbdbebf

-> Kenc  = d46c03528ecd3b1f1148e708d3e9729d
-> Kmac  = 116d156cc502da820eb57915f189ec74
\end{trace}

\begin{jcode}
package com.midnightbrewer.hardware.crypto;

import org.junit.jupiter.api.Test;
import java.util.HexFormat;
import static org.junit.jupiter.api.Assertions.assertEquals;

class DesfireSessionTest {

    @Test
    void derivesBothSessionKeys() {
        byte[] authKey = HexFormat.of()
                .parseHex("00112233445566778899aabbccddeeff");
        byte[] rndA = HexFormat.of()
                .parseHex("a0a1a2a3a4a5a6a7a8a9aaabacadaeaf");
        byte[] rndB = HexFormat.of()
                .parseHex("b0b1b2b3b4b5b6b7b8b9babbbcbdbebf");

        DesfireSession session = new DesfireSession(authKey, rndA, rndB);

        assertEquals("d46c03528ecd3b1f1148e708d3e9729d",
                HexFormat.of().formatHex(session.encKey()));
        assertEquals("116d156cc502da820eb57915f189ec74",
                HexFormat.of().formatHex(session.macKey()));
    }
}
\end{jcode}

\begin{verifybox}
\begin{shell}
mvn test -Dtest=DesfireSessionTest
\end{shell}
\textbf{If only one key is wrong}, your \code{sv[0]}/\code{sv[1]} labels are
swapped. \textbf{If both are wrong}, the SV layout is off --- print your \code{sv}
as hex and compare byte-by-byte against the diagram. The first eight bytes should
read \hex{a5 5a 00 01 00 80 a0 a1}, and eyeballing those finds most mistakes in
seconds.
\end{verifybox}

\section{Traps}

\begin{traptable}
Both keys wrong & SV layout. Print it and compare against the spec. \\
Keys swapped & Label bytes \hex{0xA5}/\hex{0x5A} reversed. \\
\code{ArrayIndexOutOfBounds} in the XOR & Slice lengths: \code{rndA[2..7]} is 6 bytes, \code{rndB[6..15]} is 10, \code{rndA[8..15]} is 8. \\
Works in the test, fails on the card & You passed \code{rndB} where \code{rndA} belongs. The order matters. \\
\end{traptable}

\section{Commit}

\begin{shell}
git add -A
git commit -m "M7b: EV2 session key derivation, verified against independent vectors"
\end{shell}
